\item Cuando la luz incide normalmente en la interfaz entre dos medios ópticos transparentes, la intensidad de la luz reflejada está dada por la expresión:
$$ S'_1 = \left( \frac{n_2 - n_1}{n_2 + n_1} \right)^2 S_1 $$
donde $S_1$ representa la intensidad de la luz incidente y $S'_1$ es la intensidad reflejada.
\begin{enumerate}
    \item ¿Qué fracción de la intensidad incidente de luz reflejada de $589\text{ nm}$ incide normalmente en una interfaz entre el aire y el vidrio corona?
    \item ¿Tiene importancia si la luz está en el aire o en el cristal cuando incide en la interfaz? Explique.
\end{enumerate}